% TEMPLATE-MILC-v3.tex - plantilla maestra UNICA de las guias MILC v3.
%
% La usan las cuatro familias de guias, cada una con su driver:
%   · Grados 6-11          -> scripts/build-guias-g11.py
%   · Bebras               -> scripts/build-guias-bebras.py
%   · Semillero            -> scripts/build-guias-semillero.py
%   · Territorio interior  -> scripts/build-guias-territorio-interior.py
%
% Antes habia tres copias casi identicas (~400 lineas iguales, 6 distintas):
% cada mejora habia que aplicarla tres veces, y en la practica no se hacia
% (el motor de recursos vivia solo en dos de ellas). Lo que varia entre
% familias esta parametrizado en 6 placeholders que cada driver llena
% (se nombran aqui SIN su sintaxis real: el builder reemplaza por texto plano,
% tambien dentro de los comentarios, y un valor multilinea rompe el preambulo):
%
%   HEADER_IZQ        encabezado izquierdo de pagina
%   HEADER_DER        encabezado derecho de pagina
%   PORTADA_ETIQUETA  rotulo sobre el numero grande (GRADO/MOMENTO/MODULO)
%   PORTADA_SERIE     linea de serie de la portada
%   PORTADA_ID        identificador de la guia en la portada
%   PORTADA_CIRCULO   el circulito de grado (°), o vacio si no aplica
%   PDF_TITULO        titulo en texto plano (sin \textbf ni comillas TeX) para
%                     los metadatos del PDF (/Title); lo produce lib_xelatex.texto_pdf
%
% Composicion (auditoria 2026-09-03): las bandas de estacion piden espacio
% (\Needspace*) y prohiben el salto justo despues (\nopagebreak), asi no quedan
% huerfanas al pie; las cajas largas (softbox, drawbox, infoband, puentebox) son
% `breakable`, asi una caja de media pagina ya no arrastra todo a la siguiente
% dejando la anterior al 40 %. Todo texto sobre mostaza o verde va en milcNegro
% (blanco sobre #E5B400 da 1,93:1 y sobre #58A923 2,95:1; ninguno pasa WCAG AA).
% El resto de claves las documenta content/guias/_SCHEMA.md. El script valida
% que no queden placeholders sin reemplazar antes de escribir el archivo final.

% Etiquetado PDF (latex-lab): /Lang, estructura y texto alternativo de las
% figuras, para lectores de pantalla. Debe ir ANTES de \documentclass.
\DocumentMetadata{lang=es-CO,tagging=on}
\documentclass[11pt,letterpaper]{article}
% headheight=24pt: la cabecera derecha lleva dos lineas (identidad de la guia
% y estacion actual). headsep se acorta en la misma medida para que el bloque
% de texto quede exactamente donde estaba con headheight=12pt.
\usepackage[margin=2.0cm,headheight=24pt,headsep=13pt]{geometry}
\usepackage[table]{xcolor}
\usepackage{fontspec}
\usepackage[spanish]{babel}
\usepackage{tikz}
% Librerías TikZ para los diagramas de las guías (bloque `recursos:`):
%  · babel  → babel en español vuelve activos caracteres como «>», y TikZ los usa
%             en sus opciones (>=stealth, ->). Sin esta librería, cualquier
%             diagrama con flechas revienta con «\language@active@arg> has an
%             extra }». Debe ir ANTES de que se dibuje nada.
%  · positioning        → sintaxis «below left=2pt and 3pt».
%  · decorations.pathreplacing → llaves ({brace}) para señalar rangos.
\usetikzlibrary{babel,positioning,decorations.pathreplacing}
\usepackage{graphicx}
% Circuitos: el trabajo de electrónica se hace hoy con micro:bit y sus sensores
% integrados (MakeCode, sin cableado), así que no se necesitan esquemáticos.
% Si algún día entran montajes con protoboard, basta descomentar esta línea:
% los diagramas .tex/.tikz declarados en `recursos:` ya se incrustan solos.
% \usepackage{circuitikz}
\usepackage[most]{tcolorbox}
\usepackage{tabularx}
\usepackage{array}
\usepackage{fancyhdr}
\usepackage{titlesec}
\usepackage[document]{ragged2e}  % bandera a la derecha en todo el cuerpo
\usepackage{enumitem}
\usepackage{microtype}
\usepackage{etoolbox}
\usepackage{needspace}
% hyperref al final del preambulo: metadatos del PDF (titulo, autor, idioma,
% familia), marcadores por estacion (\pdfbookmark) y enlaces sin recuadro.
\usepackage[hidelinks,
  pdftitle={Sustentar el mini estudio — cuatro minutos y las preguntas de los demás},
  pdfauthor={PhD. Álvaro Cárdenas Orozco},
  pdfsubject={Tecnología e Informática · Grado 8 · Guía 10 · MILC v3},
  pdflang={es-CO},
  bookmarksopen=true,bookmarksnumbered=false]{hyperref}

% Helvetica Neue no trae la flecha U+2192, asi que cada «→» escrito literalmente
% en el YAML se compilaba a NADA, en silencio (462 flechas en 61 guias: rutas del
% tipo «paleta → tipografia → lockup» salian rotas). Mapeamos el caracter a su
% version matematica, que si tiene glifo. Asi el autor puede seguir escribiendo
% «→» comodamente en el YAML.
\usepackage{newunicodechar}
\newunicodechar{→}{\ensuremath{\rightarrow}}
% Mismo problema con los simbolos de marca y vinetas: el «✓» de «una palomita ✓»
% salia como cuadrito vacio en el PDF de todas las guias que lo usaban.
\usepackage{pifont}
\newunicodechar{✓}{\ding{51}}
\newunicodechar{✗}{\ding{55}}
\newunicodechar{✔}{\ding{52}}
\newunicodechar{•}{\textbullet}
\newunicodechar{≥}{\ensuremath{\geq}}
\newunicodechar{≤}{\ensuremath{\leq}}

\setmainfont{Helvetica Neue}
\setsansfont{Helvetica Neue}
\setlength{\parindent}{0pt}
\setlength{\parskip}{6pt}
% Sin particion de palabras: en 6.º-7.º una palabra cortada por guion al final
% de linea («Comuni-cacion») frena la lectura mucho mas que una linea corta.
\hyphenpenalty=10000
\exhyphenpenalty=10000
\pretolerance=2000
\tolerance=3000
\setlength{\emergencystretch}{3em}
\linespread{1.10}
\renewcommand{\arraystretch}{1.25}
\setlist{leftmargin=5mm,itemsep=1mm,topsep=1mm}
\newcolumntype{Y}{>{\RaggedRight\arraybackslash}X}

\definecolor{milcMagenta}{HTML}{E6007E}
\definecolor{milcVino}{HTML}{5A0038}
\definecolor{milcTurquesa}{HTML}{0093A5}
\definecolor{milcVerde}{HTML}{58A923}
\definecolor{milcMostaza}{HTML}{E5B400}
\definecolor{milcOcre}{HTML}{B86600}
\definecolor{milcNegro}{HTML}{241816}
\definecolor{milcGris}{HTML}{F7F7F4}
\definecolor{milcLinea}{HTML}{E8E2DD}
\definecolor{milcGradePrimary}{HTML}{FF6600}
\definecolor{milcGradeDark}{HTML}{9A3E00}
\definecolor{milcGradeSoft}{HTML}{FFE4D0}
\definecolor{milcGradeAccent}{HTML}{CFE7FF}
\definecolor{milcGradeText}{HTML}{FFFFFF}
\color{milcNegro}

\pagestyle{fancy}
\fancyhf{}
% Cabecera de dos lineas: identidad de la guia arriba y, a la derecha y en
% gris, la estacion en curso (\markright desde \milcestacion). Antes de la
% primera estacion la segunda linea va vacia; el \strut mantiene la altura
% para que la linea de arriba no baile entre paginas.
\lhead{\small Tecnología e Informática · Grado 8\\[-1pt]{\footnotesize\strut}}
\rhead{\small Guía 10 · MILC v3\\[-1pt]{\footnotesize\strut\color{milcNegro!65}\rightmark}}
\cfoot{\small\thepage}
\renewcommand{\headrulewidth}{0pt}

% Sobre mostaza (#E5B400) y verde (#58A923) el texto va en milcNegro; sobre el
% resto de la paleta, en blanco. Se usa dentro de fonttitle/fontupper, donde un
% \color posterior manda sobre coltitle.
\newcommand{\milctintasobre}[1]{%
  \ifboolexpr{test{\ifstrequal{#1}{milcMostaza}} or test{\ifstrequal{#1}{milcVerde}}}%
    {\color{milcNegro}}{\color{white}}}

\newtcolorbox{titlebox}[1]{
  enhanced,colback=#1,colframe=#1,arc=7mm,boxrule=0pt,
  left=7mm,right=7mm,top=3mm,bottom=3mm,
  before skip=8pt,after skip=8pt,
  fontupper=\sffamily\bfseries\Large\color{white}
}

\newtcolorbox{softbox}[3]{
  enhanced,breakable,pad at break=3mm,
  colback=#2,colframe=#1,borderline west={4pt}{0pt}{#1},
  arc=2mm,boxrule=.4pt,left=4mm,right=4mm,top=3mm,bottom=3mm,
  before skip=6pt,after skip=8pt,title={#3},coltitle=white,
  colbacktitle=#1,fonttitle=\sffamily\bfseries\milctintasobre{#1},fontupper=\RaggedRight,
  attach boxed title to top left={xshift=3mm,yshift=-2mm},
  boxed title style={arc=2mm, boxrule=0pt}
}

\newtcolorbox{drawbox}[2]{
  enhanced,breakable,pad at break=4mm,
  colback=white,colframe=#1,arc=4mm,boxrule=1pt,
  left=5mm,right=5mm,top=4mm,bottom=4mm,before skip=8pt,after skip=8pt,
  title={#2},coltitle=white,colbacktitle=#1,fonttitle=\sffamily\bfseries\milctintasobre{#1},
  fontupper=\RaggedRight,
  attach boxed title to top left={xshift=4mm,yshift=-2mm},
  boxed title style={arc=3mm, boxrule=0pt}
}

\newtcolorbox{infoband}[1]{
  enhanced,breakable,pad at break=2mm,colback=#1!8,colframe=#1!50,arc=2mm,boxrule=.5pt,
  left=4mm,right=4mm,top=2mm,bottom=2mm,before skip=4pt,after skip=4pt,
  fontupper=\small\RaggedRight
}

% Puente narrativo entre secciones (contrato editorial Conéctate).
% Da continuidad: "Acabas de X. Ahora vas a Y porque Z."
% Visualmente: banda izquierda en color del grado, texto en itálica pequeña.
\newtcolorbox{puentebox}{
  enhanced,breakable,pad at break=2mm,colback=milcGradeSoft,colframe=milcGradeDark,
  borderline west={3pt}{0pt}{milcGradeDark},
  arc=0mm,boxrule=0pt,left=5mm,right=4mm,top=2.5mm,bottom=2.5mm,
  before skip=4pt,after skip=8pt,
  fontupper=\itshape\small\color{milcNegro}\RaggedRight
}
% La flecha va en modo matematico a proposito: Helvetica Neue NO tiene el glifo
% U+2192, asi que \textrightarrow se compilaba a NADA («Missing character: There
% is no → in font Helvetica Neue») y el puente salia sin flecha. En modo
% matematico la flecha viene de la fuente matematica y si se dibuja.
\newcommand{\puente}[1]{%
  \begin{puentebox}%
    \textcolor{milcGradeDark}{$\rightarrow$}\hspace{2mm}#1%
  \end{puentebox}%
}

\newcommand{\linead}[1]{\noindent\color{milcLinea}\rule{#1}{0.5pt}\color{milcNegro}}
\newcommand{\respuesta}[1]{\par\vspace{2mm}\foreach \n in {1,...,#1}{\noindent\color{milcLinea}\rule{\linewidth}{0.5pt}\par\vspace{5mm}}\color{milcNegro}}
\newcommand{\checkbox}{\textcolor{milcGradeDark}{\fbox{\phantom{x}}}\hspace{5pt}}

% ─── Listas aireadas para las guías socioemocionales ────────────────────
% Componer las verificaciones como «\checkbox texto\\» deja el texto que salta
% de línea debajo del cuadrito y apelmaza el bloque. Estos entornos alinean la
% continuación y airean los ítems. Los usa Territorio Interior; cualquier
% familia puede hacerlo.
\newlist{checklista}{itemize}{1}
\setlist[checklista]{
  label=\textcolor{milcGradeDark}{\fbox{\phantom{x}}},
  leftmargin=9mm, labelsep=3.2mm, itemsep=2.4mm,
  topsep=2.5mm, parsep=0pt, partopsep=0pt, before=\RaggedRight
}
\newlist{pasoslista}{enumerate}{1}
\setlist[pasoslista]{
  label=\textcolor{milcGradeDark}{\sffamily\bfseries\arabic*.},
  leftmargin=8mm, labelsep=3mm, itemsep=2.8mm,
  topsep=1mm, parsep=0pt, partopsep=0pt, before=\RaggedRight
}

\newcommand{\phasecard}[4]{
\begin{tcolorbox}[enhanced,colback=#3,colframe=#2,arc=3mm,boxrule=.5pt,height=3.05cm,valign=center,halign=center,left=1.5mm,right=1.5mm,top=2mm,bottom=2mm]
{\sffamily\bfseries\fontsize{22}{24}\selectfont\textcolor{#2}{#1}}\par
{\sffamily\bfseries\small #4}
\end{tcolorbox}
}

% ── Piezas del rediseno 2026 (legibilidad 6.º-11.º) ───────────────────
% Banda de estacion: reemplaza el titlebox macizo. Titulo a la izquierda,
% dato practico (tiempo/modalidad) a la derecha, en una sola linea.
% Nunca huerfana: pide 9 lineas libres (o salta de pagina rellenando con
% \vfill) y prohibe el salto justo despues de la banda. Nueve y no seis:
% la caja partible que sigue necesita ~80pt (titulo adosado, relleno y dos
% lineas) para arrancar en la misma pagina; con menos, tcolorbox la manda
% entera a la siguiente y la banda se queda sola. El penalty 10000 va
% en `after`, ANTES del espacio posterior: un `after skip` (aunque sea 0pt)
% es un \vskip tras la caja, y ese glue es un punto de corte legal que un
% \nopagebreak escrito despues de \end{tcolorbox} ya no cierra. Cada banda
% deja marcador PDF y actualiza la cabecera derecha.
\newcounter{milcestacion}
\newcommand{\milcestacion}[2]{%
\Needspace*{9\baselineskip}%
\stepcounter{milcestacion}%
\pdfbookmark[1]{#2}{milcestacion\themilcestacion}%
\markright{#2}%
\begin{tcolorbox}[enhanced,colback=#1,colframe=#1,arc=2mm,boxrule=0pt,
  left=5mm,right=5mm,top=2.6mm,bottom=2.6mm,before skip=11pt,
  after={\par\nopagebreak\vskip7pt}]
{\sffamily\bfseries\fontsize{15}{18}\selectfont\milctintasobre{#1}#2}
\end{tcolorbox}}

% Tarjeta de paso numerada: sustituye las tablas «Pilar / Aplicacion», que
% eran formato de consulta docente, no de estudiante. El numero va en un
% circulo sobre el borde para que la secuencia se lea de un vistazo.
\newcommand{\milcpaso}[3]{%
\Needspace{4\baselineskip}%
\begin{tcolorbox}[enhanced,colback=white,colframe=milcLinea,arc=1.5mm,boxrule=.9pt,
  left=14mm,right=4mm,top=3mm,bottom=3mm,before skip=4pt,after skip=4pt,
  fontupper=\RaggedRight,
  overlay={\node[anchor=north west,circle,fill=#2,text=white,inner sep=0pt,
    minimum size=7.4mm,font=\sffamily\bfseries\small]
    at ([xshift=3.6mm,yshift=-3.2mm]frame.north west) {#1};}]
#3
\end{tcolorbox}}

% Casilla marcable de tamano util con lapiz (la anterior era \fbox{\phantom{x}},
% de alto variable segun el texto que la seguia).
\newcommand{\milcmarca}{\framebox[3.8mm]{\rule{0pt}{3.4mm}}}

% Fila marcable de lista suelta (checks de la estacion 1).
\newcommand{\milccheckrow}[1]{%
\par\vspace{1.8mm}\noindent
\makebox[6.4mm][l]{\raisebox{-.55mm}{\textcolor{milcNegro}{\milcmarca}}}%
\parbox[t]{\dimexpr\linewidth-6.4mm\relax}{\RaggedRight #1}\par}

% Celda marcable dentro de la rubrica (tabla de autoevaluacion). Misma
% sangria francesa, pero pensada para vivir dentro de una columna X.
\newcommand{\milcopcion}[1]{%
\makebox[5.2mm][l]{\raisebox{-.55mm}{\textcolor{milcNegro}{\milcmarca}}}%
\parbox[t]{\dimexpr\linewidth-5.2mm\relax}{\RaggedRight #1}}

% ── Recursos visuales de la guía ──────────────────────────────────────
% Declarados en el bloque `recursos:` del YAML y emitidos por el builder.
% \guiaFigura   → imagen rasterizada o PDF (foto, carta celeste, montaje).
% \guiaDiagrama → fuente TikZ/circuitikz (.tex) incrustada como vector:
%                 es la vía para circuitos y esquemáticos editables.
% [#1] = texto alternativo (alt) · #2 = ruta relativa al .tex · #3 = pie
% El alt llega al PDF etiquetado (/Alt) y lo lee el lector de pantalla.
\newcommand{\guiaFigura}[3][]{%
\begin{center}
\includegraphics[width=0.88\linewidth,height=0.38\textheight,keepaspectratio,alt={#1}]{#2}\\[1.5mm]
{\footnotesize\itshape\textcolor{milcNegro!75}{#3}}
\end{center}
}

\newcommand{\guiaDiagrama}[2]{%
\begin{center}
\input{#1}\\[1.5mm]
{\footnotesize\itshape\textcolor{milcNegro!75}{#2}}
\end{center}
}

\begin{document}
\thispagestyle{empty}

% ── Portada ───────────────────────────────────────────────────────────
% Antes: un rectangulo de color a pagina completa. Se veia bien en pantalla,
% pero en la fotocopiadora del colegio se comia el toner y salia gris sucio.
% Ahora la banda ocupa el tercio superior y el resto va en blanco.
\begin{tikzpicture}[remember picture,overlay]
  \path[fill=milcGradePrimary]
    (current page.north west) rectangle ([yshift=-10.6cm]current page.north east);

  \node[anchor=north west,text=milcGradeText] at ([xshift=2.0cm,yshift=-1.55cm]current page.north west)
    {\sffamily\bfseries\fontsize{12}{14}\selectfont GRADO};
  \node[anchor=north west,text=milcGradeText,opacity=.85] at ([xshift=2.0cm,yshift=-2.35cm]current page.north west)
    {\sffamily\bfseries\fontsize{8.5}{10}\selectfont SERIE GUÍAS · TECNOLOGÍA E INFORMÁTICA};
  \node[anchor=north east,text=milcGradeText,opacity=.85,align=right] at ([xshift=-2.0cm,yshift=-1.55cm]current page.north east)
    {\sffamily\bfseries\fontsize{9}{12}\selectfont I.E. Sor María Juliana\\Cartago, Valle del Cauca};

  \node[anchor=west,text=milcGradeText] (gradonum) at ([xshift=1.85cm,yshift=-7.1cm]current page.north west)
    {\sffamily\bfseries\fontsize{112}{112}\selectfont 8};
  % El circulito de grado (°) solo aplica a las guias de grado («11°»); en
  % Bebras y Semillero el numero grande es un momento/modulo, no un grado.
  % Se posiciona relativo al nodo `gradonum`, no en coordenadas absolutas.
  \draw[milcGradeText,line width=1.6pt]
    ([xshift=2.0mm,yshift=-4.5mm]gradonum.north east) circle (.15cm);

  \node[anchor=west,text=milcGradeText,text width=11.4cm,align=left] at ([xshift=1.5cm,yshift=0.5cm]gradonum.east)
    {
     {\sffamily\bfseries\fontsize{9}{11}\selectfont\MakeUppercase{Período 1 · Análisis de datos con phronesis}}\\[3mm]
     {\sffamily\bfseries\fontsize{21}{25}\selectfont\setlength{\baselineskip}{27pt}Sustentar en\\[4pt]cuatro minutos}\\[3.5mm]
     {\sffamily\bfseries\fontsize{15}{18}\selectfont Guía 10-8-TIC}
    };
\end{tikzpicture}

\vspace*{9.4cm}
\vfill

{\sffamily\bfseries\fontsize{8}{10}\selectfont\textcolor{milcGradeDark}{LO QUE TE LLEVAS HOY}}

\begin{tcolorbox}[enhanced,colback=white,colframe=milcNegro,arc=2mm,boxrule=1.6pt,
  left=5mm,right=5mm,top=4mm,bottom=4mm,before skip=3pt,after skip=12pt,
  fontupper=\RaggedRight\fontsize{11}{15.5}\selectfont]
Una sustentación oral de cuatro minutos de tu mini estudio, con cinco diapositivas (pregunta, método, hallazgo, gráfico, decisión y limitaciones). Tres a cinco preguntas del grupo respondidas con honestidad. Y una autoevaluación con una fortaleza y una mejora concreta para la próxima vez.
\end{tcolorbox}

\begin{tcolorbox}[enhanced,colback=milcGris,colframe=milcGris,arc=2mm,boxrule=0pt,
  left=5mm,right=5mm,top=3.5mm,bottom=3.5mm,after skip=12pt]
{\sffamily\bfseries\fontsize{8}{10}\selectfont\textcolor{milcNegro!55}{ESTA GUÍA ES DE}}\\[3mm]
\begin{tabularx}{\linewidth}{X p{3.2cm} p{3.2cm}}
\linead{\linewidth} & \linead{3.0cm} & \linead{3.0cm}\\[-1mm]
{\fontsize{8.5}{10}\selectfont\textcolor{milcNegro!55}{Nombre completo}} &
{\fontsize{8.5}{10}\selectfont\textcolor{milcNegro!55}{Grupo}} &
{\fontsize{8.5}{10}\selectfont\textcolor{milcNegro!55}{Fecha}}\\
\end{tabularx}
\end{tcolorbox}

\vfill

\noindent\textcolor{milcLinea}{\rule{\linewidth}{1pt}}\\[2mm]
{\sffamily\bfseries\fontsize{9.5}{12}\selectfont Autor: Dr. Álvaro Cárdenas Orozco}\\[1mm]
{\sffamily\fontsize{8.5}{11}\selectfont\textcolor{milcNegro!55}{Metodología MILC · apertura ancestral · pensamiento computacional · triángulo Dussel–estoicismo–Floridi}}

\newpage

\section*{Apertura: Sustentar el mini estudio --- cuatro minutos y las preguntas de los demás}

\begin{softbox}{milcGradeDark}{milcGradeSoft}{Saber ancestral}
Entre los misak, en el resguardo de Guambía (Silvia, Cauca), la escuela no empieza en un salón: empieza en la cocina. Llaman \textbf{nak chak} a la cocina y \textbf{nak kuk} al fogón. Los mayores dicen que ahí «se dinamiza la educación propia» (Montano Morales, 2025). Alrededor del fuego se cocina, se come, se cuentan las cosas del día y se resuelven los desacuerdos. Un asunto no se resuelve donde ocurre: se lleva al fogón, se cuenta, y los demás preguntan. Quien cuenta no expone; se somete a las preguntas. Eso es sustentar. La \textbf{cara de exclusión}: cocinar entre los misak se concentra en el mundo femenino, y el fogón fue espacio de palabra y también de reparto desigual del trabajo. Y una advertencia: el nak chak es de un pueblo que defiende su territorio; lo que hacemos hoy es mucho más pequeño y aprende de eso, no lo imita. Hoy vas a llevar tu mini estudio al fogón del salón: cuatro minutos para contarlo y el tiempo que haga falta para las preguntas.\par\smallskip{\footnotesize\textcolor{milcNegro!70}{\textbf{Fuente.} Montano Morales, J. (2025). Comer y vivir bien en el mundo indígena misak: «sentipensar el cuidado de la vida». \emph{Revista de Estudios Sociales}, (91), 163--174.}}
\end{softbox}

\begin{softbox}{milcTurquesa}{milcGris}{Saber contemporáneo}
Una \textbf{sustentación} es contar en público un trabajo propio con tres condiciones. Tiempo limitado: aquí, cuatro minutos. Un apoyo visual que ayuda sin reemplazar tu voz: cinco diapositivas. Y preguntas del público, respondidas con honestidad. El tiempo corto no es castigo: te obliga a decidir qué es esencial. Para un estudio de ocho etapas son unos treinta segundos por etapa. Lo que no cabe, o no era esencial, o todavía no lo entiendes del todo. Las preguntas tampoco son ataque: son la parte donde se ve si el estudio se sostiene. «No lo resolví, porque mi muestra fue pequeña» es una respuesta mejor que una explicación inventada.
\end{softbox}

\begin{softbox}{milcMostaza}{milcGris}{Pregunta-puente}
\textit{En el nak chak, quien cuenta el asunto no se defiende: escucha las preguntas. Cuando un compañero te pregunte «¿y a quién no encuestaste?», ¿vas a responder o a defenderte? ¿Qué diferencia hay?}
\end{softbox}

\begin{softbox}{milcVerde}{milcGris}{Saber y saber hacer}
Al terminar podrás: (1) \textbf{analizar} tu informe para reducirlo a cinco ideas que quepan en cuatro minutos; (2) \textbf{crear} cinco diapositivas que ayuden sin reemplazar tu voz; (3) \textbf{evaluar} tu sustentación y la de otros con criterios claros; (4) \textbf{explicar} con honestidad lo que tu estudio no pudo saber.
\end{softbox}



\small
\begin{tabularx}{\textwidth}{>{\bfseries}p{3.0cm}Y}
\rowcolor{milcGradePrimary!12}Autor & Dr. Álvaro Cárdenas Orozco\\
DBA & Tratamiento de datos cuantitativos con criterio ético (MEN, T\&I)\\
Referentes & Phronesis aristotélica · Floridi (ética del dato) · Estoicismo (ver lo que es)\\
Producto final & Una sustentación oral de cuatro minutos de tu mini estudio, con cinco diapositivas (pregunta, método, hallazgo, gráfico, decisión y limitaciones). Tres a cinco preguntas del grupo respondidas con honestidad. Y una autoevaluación con una fortaleza y una mejora concreta para la próxima vez.\\
\end{tabularx}
\normalsize

\puente{En la sesión 9 hiciste el mini estudio. Hoy lo cuentas en cuatro minutos y respondes preguntas. Con esto cierra el periodo; el siguiente arranca con lógica y micro:bit.}

\milcestacion{milcGradeDark}{Ruta MILC: así vamos a aprender}
\begin{tabularx}{\textwidth}{>{\centering\arraybackslash}X>{\centering\arraybackslash}X>{\centering\arraybackslash}X>{\centering\arraybackslash}X}
\phasecard{1}{milcVerde}{milcGris}{Escucha\\[-1mm]\small Observo, escucho y pregunto.} &
\phasecard{2}{milcTurquesa}{milcGris}{Entender\\[-1mm]\small Sistematizo y ordeno.} &
\phasecard{3}{milcMagenta}{milcGris}{Hacer\\[-1mm]\small Construyo y aplico.} &
\phasecard{4}{milcVino}{milcGris}{Cierre\\[-1mm]\small Evalúo y reflexiono.}
\end{tabularx}


\puente{Antes de armar diapositivas, vas a contar tu estudio en voz alta con cronómetro, sin leer. Ahí descubres qué sobra y qué falta.}

\milcestacion{milcVerde}{Estación 1 de 3 · Escucha}

\begin{softbox}{milcVerde}{milcGris}{Escena para escuchar}
\textbf{Actividad 1 · ANALIZA --- Cuatro minutos sin leer} (15 min · individual). (1) Toma tu informe y ponlo boca abajo. (2) Pon el cronómetro en cuatro minutos y cuenta tu estudio en voz alta, como a un amigo: pregunta, hipótesis, método, hallazgo, decisión. (3) Al terminar, anota qué dejaste fuera por falta de tiempo. (4) Anota qué dijiste de más. (5) Anota qué parte sonó más floja y por qué.
\end{softbox}

{\sffamily\bfseries\fontsize{8}{10}\selectfont\textcolor{milcNegro!55}{MARCA CUANDO LO HAYAS HECHO}}
\milccheckrow{Conté el estudio completo en voz alta, sin leer, con cronómetro.}
\milccheckrow{Anoté qué dejé fuera y qué dije de más.}
\milccheckrow{Sé cuál parte sonó más floja.}

\begin{infoband}{milcVerde}


\textbf{Lo que va al cuaderno (Actividad 1 · ANALIZA).} \textbf{Título:} «Actividad 1 --- Cuatro minutos sin leer». \textbf{Formato:} tres líneas (qué dejé fuera / qué dije de más / qué sonó flojo) y el tiempo que tardé. \textbf{Extensión:} un tercio de página. \textbf{Sabes que terminaste cuando} las tres líneas están llenas y sabes qué vas a recortar antes de armar las diapositivas.
\end{infoband}

\puente{Ya sabes cuánto cabe en cuatro minutos. Ahora armas las cinco diapositivas con tu pareja, ensayas y os corregís.}

\milcestacion{milcTurquesa}{Estación 2 de 3 · Entender}

\begin{softbox}{milcTurquesa}{milcGris}{Lo que vas a entender}
\textbf{Actividad 2 · CREA --- Cinco diapositivas y un ensayo} (30 min · parejas). (1) Reduce tu informe a cinco ideas, una por diapositiva: la pregunta y por qué importa; el método; el hallazgo con sus cifras; el gráfico honesto; la decisión y las limitaciones. (2) Arma las cinco en PowerPoint, Canva o Slides, con máximo veinticinco palabras cada una. (3) Escribe el guion hablado: unas seiscientas palabras son cuatro minutos. (4) Ensaya frente a tu pareja con cronómetro. (5) Tu pareja te dice una cosa que no entendió y una que sobró; cambien de rol.
\end{softbox}

\milcpaso{1}{milcTurquesa}{\textbf{Las cinco diapositivas y su tiempo.} Pregunta y por qué importa: 45 segundos. Método: 45. Hallazgo con cifras: 60. Gráfico: 45. Decisión y limitaciones: 45. Cada diapositiva se lee en cinco segundos; tu voz dice el resto.}
\milcpaso{2}{milcTurquesa}{\textbf{Menos de veinticinco palabras.} Si la diapositiva tiene más, la gente la lee y deja de escucharte. Una cifra grande, un título claro, un gráfico. Lo demás lo cuentas tú.}
\milcpaso{3}{milcTurquesa}{\textbf{El guion no se lee: se recuerda.} Escríbelo para ordenarte, ensáyalo dos veces y guárdalo. Si lees, no estás sustentando. Si te pierdes, mira la diapositiva: para eso está.}
\milcpaso{4}{milcTurquesa}{\textbf{Las cifras con su contexto.} «El promedio es 3» no dice nada. «El promedio es 3 horas al día de pantalla, y el máximo fue 8» sí. Cada cifra con su unidad y con qué se compara.}

\begin{drawbox}{milcTurquesa}{La sustentación, en una mirada}
\textbf{1. Pregunta:} y por qué importa. \\[1mm] \textbf{2. Método:} a quién, cuántos, cómo. \\[1mm] \textbf{3. Hallazgo:} una frase con cifras. \\[1mm] \textbf{4. Gráfico:} honesto. \\[1mm] \textbf{5. Decisión:} y lo que no pude saber.
\end{drawbox}

\begin{softbox}{milcMostaza}{milcGris}{Errores comunes y su corrección}
(1) \textbf{Leer las diapositivas}: si lees, no sustentas. (2) \textbf{Diapositiva llena}: más de veinticinco palabras y nadie escucha. (3) \textbf{Cifra sin contexto}: «el promedio es 3,5» de qué. (4) \textbf{Inflar el hallazgo}: concluir más de lo que veinte datos aguantan. (5) \textbf{Evadir la pregunta difícil}: peor que decir «no lo resolví». (6) \textbf{Pasarse del tiempo}: es falta de ensayo, no exceso de contenido.
\end{softbox}

\begin{infoband}{milcTurquesa}


\textbf{Lo que va al cuaderno (Actividad 2 · CREA).} \textbf{Título:} «Actividad 2 --- Cinco diapositivas y un ensayo». \textbf{Formato:} las cinco ideas, una línea cada una, con su tiempo; y las dos observaciones de tu pareja. \textbf{Extensión:} media página. \textbf{Sabes que terminaste cuando} las cinco diapositivas tienen menos de veinticinco palabras, ensayaste con cronómetro y aplicaste al menos una observación de tu pareja.
\end{infoband}

\puente{Ya ensayaste. Toca sustentar frente al grupo, responder preguntas y evaluar con criterios lo que hiciste y lo que hicieron los demás.}

\milcestacion{milcMagenta}{Estación 3 de 3 · Hacer}

\begin{softbox}{milcMagenta}{milcGris}{Cómo vas a construir tu producto}
\textbf{Actividad 3 · EVALÚA --- Sustentar, responder y evaluar} (35 min · individual). (1) Sustenta frente al grupo, o en rondas de cinco, con cronómetro visible: cuatro minutos. (2) Responde de tres a cinco preguntas; si no sabes, di «no lo resolví» y por qué. (3) Mientras escuchas a otros, evalúa con los cinco criterios de la tabla y escribe una pregunta para cada uno. (4) Al final, escribe tu autoevaluación: una fortaleza concreta y una mejora concreta. (5) Entrega las diapositivas, el guion y la autoevaluación.
\end{softbox}

\milcpaso{1}{milcMagenta}{\textbf{Tu entrega tiene cuatro piezas.} (1) Las cinco diapositivas. (2) La sustentación hecha en cuatro minutos. (3) Las preguntas respondidas con honestidad. (4) La autoevaluación con una fortaleza y una mejora.}
\milcpaso{2}{milcMagenta}{\textbf{Cómo se responde una pregunta difícil.} Primero repítela con tus palabras, para asegurarte de que la entendiste. Después responde con lo que sabes. Si no lo sabes, dilo: «no lo resolví porque encuesté solo a mi salón; con más tiempo encuestaría a los tres octavos».}
\milcpaso{3}{milcMagenta}{\textbf{Cómo se evalúa a un compañero.} Con los cinco criterios, no con el gusto: ¿cumplió el tiempo?, ¿se entendió qué hizo?, ¿las cifras tenían contexto?, ¿el gráfico era honesto?, ¿dijo qué no pudo saber? Y una pregunta que le ayude, no que lo hunda.}
\milcpaso{4}{milcMagenta}{\textbf{La autoevaluación es concreta.} «Controlé el tiempo» es una fortaleza. «Estuvo bien» no. «La próxima vez ensayo tres veces» es una mejora. «Mejorar» no.}

\begin{drawbox}{milcMagenta}{Antes de entregar}
\checkbox Cinco diapositivas con menos de veinticinco palabras cada una. \\[1mm] \checkbox Guion de unas seiscientas palabras, ensayado dos veces. \\[1mm] \checkbox Sustentación dentro de los cuatro minutos. \\[1mm] \checkbox Tres a cinco preguntas respondidas; «no lo resolví» cuando aplique. \\[1mm] \checkbox Evaluación de al menos tres compañeros con los cinco criterios. \\[1mm] \checkbox Autoevaluación con una fortaleza y una mejora concretas.
\end{drawbox}

\begin{softbox}{milcMagenta}{milcGris}{Plantilla de trabajo}
\textbf{Diapositiva 1.} \textit{Pregunta y por qué importa, 45 s.} \\[1mm] \textbf{Diapositiva 2.} \textit{Método, 45 s.} \\[1mm] \textbf{Diapositiva 3.} \textit{Hallazgo con cifras, 60 s.} \\[1mm] \textbf{Diapositiva 4.} \textit{Gráfico, 45 s.} \\[1mm] \textbf{Diapositiva 5.} \textit{Decisión y limitaciones, 45 s.}
\end{softbox}

\begin{infoband}{milcMagenta}


\textbf{Lo que va al cuaderno (Actividad 3 · EVALÚA).} \textbf{Título:} «Actividad 3 --- Sustentar, responder y evaluar». \textbf{Formato:} tabla de evaluación de tres compañeros con los cinco criterios y una pregunta para cada uno; y tu autoevaluación en dos líneas. \textbf{Extensión:} una página. \textbf{Sabes que terminaste cuando} sustentaste dentro del tiempo, respondiste al menos tres preguntas, evaluaste a tres compañeros con criterios y tu mejora dice qué vas a hacer distinto.
\end{infoband}

\puente{Lo que entregas hoy: las cinco diapositivas, la sustentación hecha, las preguntas respondidas y tu autoevaluación. La prueba de calidad: quien te escuchó puede decir en una frase qué hiciste y qué encontraste.}

\milcestacion{milcMostaza}{Producto de la sesión}
\textbf{Sustentación de cuatro minutos + cinco diapositivas + preguntas respondidas + autoevaluación}.
\begin{softbox}{milcMostaza}{milcGris}{Criterios de calidad}
(1) Sustentaste dentro de los cuatro minutos, sin leer las diapositivas. (2) Las cinco diapositivas tienen menos de veinticinco palabras y siguen el orden de las cinco ideas. (3) Las cifras del hallazgo llevan unidad y comparación. (4) Respondiste de tres a cinco preguntas, con «no lo resolví» cuando aplicaba. (5) Evaluaste a tres compañeros con los cinco criterios y una pregunta útil. (6) Tu autoevaluación tiene una fortaleza y una mejora concretas.
\end{softbox}

\puente{Antes de cerrar, mira el periodo entero desde cinco lados. Diez sesiones, un estudio propio y cuatro minutos para contarlo.}

\milcestacion{milcVino}{Evaluación liberadora · cinco dimensiones}

\begin{softbox}{milcVino}{milcGris}{Sentido de la evaluación}
Evaluar no es solo revisar una nota. Es reconocer qué aprendí, qué cambió en mí, qué emociones manejé, cómo me relaciono con la ciudad y la comunidad, y qué le dejaré a quien viene detrás.
\end{softbox}

\textbf{Mi reflexión liberadora en cinco dimensiones.} Completa cada frase iniciada:

\small
\begin{tabularx}{\textwidth}{>{\bfseries\RaggedRight\arraybackslash}p{3.4cm}Y>{\RaggedRight\arraybackslash}p{4.6cm}}
\rowcolor{milcVino}\color{white}Dimensión & \color{white}\bfseries Mi respuesta (frame starter) & \color{white}\bfseries Algo que descubrí\\
1. Personal & Lo que más cambió en mí fue \linead{6.5cm} & \linead{4.2cm}\\
2. Emocional & La emoción más fuerte fue \linead{2.5cm} y la regulé \linead{3cm} & \linead{4.2cm}\\
3. Ciudadana & En democracia esto importa porque \linead{6cm} & \linead{4.2cm}\\
4. Local (Cartago) & En mi barrio sería útil para \linead{6.5cm} & \linead{4.2cm}\\
5. Intergeneracional & A mi abuela/abuelo le diría \linead{6.5cm} & \linead{4.2cm}\\
\end{tabularx}
\normalsize

\begin{infoband}{milcVino}
\textbf{Para la próxima sustentación.} Vuelve a esta tabla en seis meses. Verás que las respuestas cambian — no porque lo aprendido cambie, sino porque tú cambias. La evaluación liberadora es una conversación contigo a través del tiempo.
\end{infoband}


\puente{Para terminar el periodo, tres ideas sobre escuchar y contar, con palabras de esta guía: Dussel sobre el respeto como silencio del que escucha, Marco Aurelio sobre atender a lo que dice el otro, Floridi sobre la atención como un bien escaso.}

\milcestacion{milcGradeDark}{Tres ideas para cerrar · Dussel, un estoico y Floridi}

\begin{softbox}{milcVino}{milcGris}{Enrique Dussel · lente del nosotros}
\textit{El respeto es silencio; no el del que nada tiene que decir, sino el del que tiene todo que escuchar.}

{\footnotesize\itshape\color{milcNegro!70}--- idea tomada de Enrique Dussel, contada con palabras de esta guía. No es una cita textual.}

En el nak chak, quien cuenta calla mientras los demás preguntan. Cuando un compañero sustenta, tu silencio atento es lo que hace posible que su estudio se sostenga o se corrija. Y cuando te preguntan a ti, escuchar la pregunta completa antes de responder es respeto.

\textbf{¿Escuché la pregunta entera, o empecé a responder antes de que terminara?}
\end{softbox}
\linead{\linewidth}\\[1mm]
\linead{\linewidth}

\begin{softbox}{milcMostaza}{milcGris}{Estoicismo · Marco Aurelio · cuidado interior}
\textit{Acostúmbrate a atender a lo que dice el otro, y a meterte, hasta donde puedas, en lo que piensa.}

{\footnotesize\itshape\color{milcNegro!70}--- idea tomada de Marco Aurelio, contada con palabras de esta guía. No es una cita textual.}

Una pregunta difícil no es un ataque: es alguien que pensó en tu estudio. Repetirla con tus palabras antes de responder es meterte en lo que piensa. De ahí salen las mejores respuestas, y las mejores mejoras.

\textbf{¿Qué pregunta de hoy me mostró algo de mi estudio que yo no había visto?}
\end{softbox}
\linead{\linewidth}\\[1mm]
\linead{\linewidth}

\begin{softbox}{milcTurquesa}{milcGris}{Luciano Floridi y la Onlife Initiative · lente de la infoesfera}
\textit{La atención de las personas es un bien finito, precioso y escaso.}

{\footnotesize\itshape\color{milcNegro!70}--- idea tomada de Luciano Floridi y la Onlife Initiative, contada con palabras de esta guía. No es una cita textual.}

Cuatro minutos son cuatro minutos de la atención de treinta personas. Una diapositiva llena la gasta; una cifra con contexto la aprovecha. Sustentar bien es cuidar la atención de los demás, no llenarla.

\textbf{¿Qué de mi sustentación gastó la atención del grupo sin devolverle nada?}
\end{softbox}
\linead{\linewidth}\\[1mm]
\linead{\linewidth}

\begin{infoband}{milcNegro}
\textbf{Nota para el docente.} Las tres ideas están contadas con palabras de esta guía a partir del pensamiento de cada autor; no son citas textuales. De 9.º en adelante se trabajan con la obra y su referencia.
\end{infoband}

\milcestacion{milcVino}{Mi autoevaluación · marca una casilla por fila}

{\small
\setlength{\extrarowheight}{2.2pt}
\begin{tabularx}{\textwidth}{>{\RaggedRight\arraybackslash\bfseries}p{3.0cm}YYY}
\rowcolor{milcVino}\color{white}\bfseries Criterio & \color{white}\bfseries Lo logré & \color{white}\bfseries En proceso & \color{white}\bfseries Necesito apoyo\\[.6mm]
Comprensión del tema & \milcopcion{Explico las ideas con mis palabras.} & \milcopcion{Reconozco las ideas, falta precisión.} & \milcopcion{Necesito repasar lo básico.}\\[.8mm]
\rowcolor{milcGris}Pensamiento computacional & \milcopcion{Usé los pilares al armar mi producto.} & \milcopcion{Usé algunos pilares.} & \milcopcion{Necesito releer la estación 2.}\\[.8mm]
Aplicación y producto & \milcopcion{Mi producto cumple los criterios.} & \milcopcion{Está armado, con ajustes pendientes.} & \milcopcion{Aún está en construcción.}\\[.8mm]
\rowcolor{milcGris}Reflexión 5 dimensiones & \milcopcion{Respondí cada una con honestidad.} & \milcopcion{Respondí algunas con detalle.} & \milcopcion{Mis respuestas son muy generales.}\\[.8mm]
Triángulo de pensamiento & \milcopcion{Conecté las 3 voces con mi trabajo.} & \milcopcion{Conecté al menos una.} & \milcopcion{No las relaciono con mi proceso.}\\[.8mm]
\end{tabularx}}

\vspace{3mm}
\textbf{Hoy entendí que:}\\[1mm]
\linead{\linewidth}\\[1mm]
\linead{\linewidth}

\textbf{Mi compromiso para la próxima sesión es:}\\[1mm]
\linead{\linewidth}\\[1mm]
\linead{\linewidth}

\begin{softbox}{milcMostaza}{milcGris}{Cita de despedida · Marco Aurelio}
\textit{``El obstáculo en el camino se vuelve el camino. Lo que estorba al camino se vuelve camino.''}\\[1mm]
Cada error de hoy, bien observado, se convierte en pieza del camino que pisarás mañana.
\end{softbox}

\small
\begin{tabularx}{\textwidth}{>{\bfseries}p{3.6cm}Y}
\rowcolor{milcGradeDark!12}Nombre completo: & \linead{10cm}\\
Fecha: & \linead{4cm} \quad Curso/grupo: \linead{4cm}\\
Mi firma: & \linead{9cm}\\
\end{tabularx}
\normalsize

\begin{infoband}{milcGradeDark}
\textbf{Cierre del periodo.} Guarda tu mini estudio, las diapositivas y la autoevaluación: son tu evidencia del periodo 1, «Datos con phronesis». El periodo 2 arranca con lógica y micro:bit: condiciones, AND, OR y NOT. La mejora que escribiste hoy es lo primero que vas a probar en la próxima sustentación.
\end{infoband}

\end{document}
